\documentclass[spanish,letterpaper,oneside]{article}

% Datos del documento
\def\documenttitle {Conocer, hacer, ser y vivir juntos en la enseñanza técnica}
\def\documentsubtitle {Fundamentos para la enseñanza y el aprendizaje I}
\def\documentsubject {Actividad 4 · Tipos de aprendizaje y aplicación docente}
\def\documentauthor {Mtro. Martín Jonathan de la Cruz Muñoz}
\def\coursename {Fundamentos para la enseñanza y el aprendizaje I}
\def\coursecode {MAEM}

% Identidad institucional
\def\universityname {Instituto de Investigación, Innovación y Estudios de Posgrado para la Educación del Estado de Nuevo León}
\def\universityfaculty {Maestría en Aprendizaje y Enseñanza de las Matemáticas}
\def\universitydepartment {Fundamentos para la enseñanza y el aprendizaje I}
\def\universitydepartmentimage {img/iiiepe/marca-iiiepe.png}
\def\universitydepartmentimagecfg {height=1.75cm}
\def\universitylocation {Monterrey, Nuevo León}
\def\studentfullname {Mtro. Martín Jonathan de la Cruz Muñoz}
\def\studentid {00874}
\def\studentemail {martin.delacruz@campus.iiiepe.edu.mx}
\def\programperiod {1er tetramestre, marzo--julio de 2026}
\def\programmodality {Modalidad presencial}
\def\programshort {MAEM}

\def\authortable {
  \begin{tabular}{ll}
    \textbf{Alumno:} & \studentfullname \\
    Matrícula: & \studentid \\
    Correo institucional: & \studentemail \\
    Unidad: & \coursename \\
    Programa: & \programshort \\
    Periodo: & \programperiod \\
    & \\
    \multicolumn{2}{l}{Fecha de realización: 13 de octubre de 2026} \\
    \multicolumn{2}{l}{\universitylocation}
  \end{tabular}
}

\input{template}

\usepackage{helvet}
\renewcommand{\familydefault}{\sfdefault}
\usepackage{array}
\usepackage{longtable}
\usepackage{pdflscape}
\usepackage{tikz}
\usepackage[most]{tcolorbox}
\usetikzlibrary{arrows.meta,positioning,calc,fit,backgrounds}

\definecolor{iiiepeNavy}{HTML}{163B5C}
\definecolor{iiiepeBlue}{HTML}{1F6F9F}
\definecolor{iiiepeAqua}{HTML}{20B7A8}
\definecolor{iiiepeOrange}{HTML}{F28C28}
\definecolor{iiiepePaper}{HTML}{F6F8F7}
\definecolor{iiiepeInk}{HTML}{1E2A32}
\def\portraittitlecolor {iiiepeNavy}
\setcitestyle{authoryear,open={(},close={)}}

\newcommand{\insertcoverwatermark}{%
  \AddToShipoutPictureBG*{%
    \begin{tikzpicture}[remember picture,overlay]
      \node[opacity=.09,inner sep=0pt] at (current page.center)
      {\includegraphics[width=.58\paperwidth]{img/iiiepe/logo-puntos.png}};
    \end{tikzpicture}%
  }%
}

\newtcolorbox{reportelecturabox}[1][]{%
  enhanced,breakable,
  colback=iiiepePaper,colframe=iiiepeNavy,
  boxrule=1pt,leftrule=4pt,arc=2pt,
  left=8pt,right=8pt,top=6pt,bottom=6pt,
  before skip=6pt,after skip=10pt,
  fonttitle=\bfseries\color{white},coltitle=white,
  colbacktitle=iiiepeNavy,
  attach boxed title to top left={xshift=8pt,yshift=-2pt},
  boxed title style={colback=iiiepeNavy,colframe=iiiepeNavy,arc=1pt},
  title={Reporte de lectura},#1}

\newtcolorbox{fichalecturabox}[1][]{%
  enhanced,breakable,
  colback=iiiepeAqua!7,colframe=iiiepeAqua!70!iiiepeNavy,
  boxrule=.6pt,arc=1pt,
  left=7pt,right=7pt,top=4pt,bottom=4pt,
  fontupper=\small,before upper={\setstretch{1.0}},
  fonttitle=\bfseries\color{iiiepeNavy},coltitle=iiiepeNavy,
  colbacktitle=iiiepeAqua!7,
  title={Ficha de la lectura},#1}



\usepackage{tabularx,colortbl,pgfplots,enumitem,needspace}\pgfplotsset{compat=1.18}\definecolor{iiiepePale}{HTML}{EEF5F8}
\setlength{\parskip}{5pt}\setlength{\parindent}{0pt}
\begin{document}\insertcoverwatermark\templatePortrait\templatePagecfg\templateFinalcfg\fancyfoot{}\fancyfoot[L]{\scriptsize IIIEPE · MAEM}\fancyfoot[C]{\scriptsize\thepage}\fancyfoot[R]{\scriptsize Actividad 4}\setstretch{1.3}\raggedbottom
\section{Cuatro aprendizajes que deben articularse}
La lectura de Scott (2015b) organiza la reflexión sobre lo que se necesita aprender en el siglo XXI a partir de conocer, hacer, ser y vivir juntos. Estos pilares no describen estilos fijos de estudiantes ni cuatro tipos de personas. Son dimensiones complementarias de la formación. En mis clases técnicas y de ingeniería, permiten revisar si una tarea desarrolla comprensión, actuación, autonomía y convivencia, en lugar de valorar solo su ejecución.


{\small\setstretch{1.1}\renewcommand{\arraystretch}{1.25}
\begin{longtable}{|>{\raggedright\arraybackslash}p{\dimexpr0.22\linewidth-2\tabcolsep-2\arrayrulewidth\relax}|>{\raggedright\arraybackslash}p{\dimexpr0.78\linewidth-2\tabcolsep-2\arrayrulewidth\relax}|}\hline
\rowcolor{iiiepePale} \textbf{Pilar} & \textbf{Sentido y aplicación propuesta} \\\hline
\endfirsthead
\hline
\rowcolor{iiiepePale} \textbf{Pilar} & \textbf{Sentido y aplicación propuesta} \\\hline
\endhead
Aprender a conocer & Construir comprensión, formular preguntas y examinar información. En una calibración, identificar variables, unidades y supuestos; comparar un modelo proporcional con uno que tiene término inicial. La evidencia sería una explicación de por qué se elige el modelo. \\\hline
Aprender a hacer & Utilizar conocimientos para resolver y revisar un problema. Diseñar un procedimiento de medición, calcular un error y ajustar una decisión. La evidencia sería un registro trazable, con cálculos y justificación del cambio. \\\hline
\end{longtable}}



{\small\setstretch{1.1}\renewcommand{\arraystretch}{1.25}
\begin{longtable}{|>{\raggedright\arraybackslash}p{\dimexpr0.22\linewidth-2\tabcolsep-2\arrayrulewidth\relax}|>{\raggedright\arraybackslash}p{\dimexpr0.78\linewidth-2\tabcolsep-2\arrayrulewidth\relax}|}\hline
\rowcolor{iiiepePale} \textbf{Pilar} & \textbf{Sentido y aplicación propuesta} \\\hline
\endfirsthead
\hline
\rowcolor{iiiepePale} \textbf{Pilar} & \textbf{Sentido y aplicación propuesta} \\\hline
\endhead
Aprender a ser & Desarrollar autonomía, responsabilidad y juicio. Reconocer una incertidumbre, pedir apoyo y decidir con criterios de seguridad y respeto. La evidencia sería una reflexión sustentada sobre una decisión, sin confundir autonomía con trabajar sin acompañamiento. \\\hline
Aprender a vivir juntos & Participar, escuchar y construir acuerdos con otras personas. Rotar roles técnicos y argumentativos, contrastar procedimientos y ofrecer una devolución respetuosa. La evidencia sería la distribución de decisiones y la revisión compartida de una solución. \\\hline
\end{longtable}}


\section{Aplicación integrada en una secuencia}
Propondría comparar dos estrategias de mantenimiento mediante un conjunto de datos didácticos. En el inicio, cada estudiante anticiparía qué significa ``mejor'' y qué información necesita. En el desarrollo, los equipos calcularían promedios y dispersión, identificarían límites del conjunto y prepararían una recomendación. En el cierre, cada integrante explicaría una decisión y revisaría su primera postura.
Aprender a conocer aparecería al comprender las medidas estadísticas; aprender a hacer, al calcular y representar; aprender a ser, al reconocer límites y responsabilidad; aprender a vivir juntos, al construir y revisar el argumento. La actividad no se dividiría en cuatro tareas independientes: una misma decisión exige las cuatro dimensiones.
\section{Evaluación y apoyos}
Evaluaría comprensión conceptual, procedimiento, justificación y colaboración mediante evidencias diferenciadas. Quien tenga dificultades de cálculo podría utilizar una tabla guiada y conservaría la responsabilidad de interpretar; quien avance con mayor autonomía podría analizar un valor extremo y explicar su efecto. El apoyo no elimina la exigencia de pensar.
La retroalimentación incluiría una acción realizable: revisar una unidad, contrastar dos resultados o justificar un criterio. Pediría una segunda versión para que la evaluación participe en el aprendizaje. La valoración de actitudes requiere observación real; no debe inferirse del formato del documento.
\section{Un problema que integra los cuatro pilares}
Propongo comparar dos planes de mantenimiento: uno cobra $250+80h$ pesos y otro $100+110h$, para $h\ge0$. El punto de igualdad es $h=5$ y el costo común es 650 pesos. Antes de cinco horas conviene el segundo plan; después, el primero, suponiendo servicios equivalentes y sin cargos adicionales.
Aprender a conocer exige interpretar pendiente y cuota fija; aprender a hacer, construir tabla, gráfica y solución algebraica; aprender a ser, reconocer límites y sostener una recomendación responsable; aprender a vivir juntos, contrastar razones y acordar una comunicación compartida. Scott (2015b, pp. 4--8) vincula los pilares con una formación que supera la reproducción de información; mi aplicación concreta es organizar una decisión que los necesite conjuntamente.

{\small\setstretch{1.1}\renewcommand{\arraystretch}{1.25}
\begin{longtable}{|>{\raggedright\arraybackslash}p{\dimexpr0.32\linewidth-2\tabcolsep-2\arrayrulewidth\relax}|>{\raggedright\arraybackslash}p{\dimexpr0.68\linewidth-2\tabcolsep-2\arrayrulewidth\relax}|}\hline
\rowcolor{iiiepePale} \textbf{Evidencia} & \textbf{Pregunta para revisar comprensión} \\\hline
\endfirsthead
\hline
\rowcolor{iiiepePale} \textbf{Evidencia} & \textbf{Pregunta para revisar comprensión} \\\hline
\endhead
Tabla y ecuación & ¿Qué representa el costo cuando $h=0$ y por qué las dos rectas se cruzan? \\\hline
Decisión argumentada & ¿Cambiaría tu recomendación si un servicio incluye refacciones y el otro no? \\\hline
Revisión individual & ¿Qué idea modificaste después de escuchar otra estrategia? \\\hline
\end{longtable}}

\section{Conclusión reflexiva}
Mi punto de partida ha sido organizar conocimiento y trabajo; estos pilares me ayudan a ampliar la finalidad de esa organización. En preparatoria y licenciatura, formar para resolver problemas debe incluir la capacidad de comprender, actuar con responsabilidad y convivir. Aplicaría esta mirada al diseñar tareas donde cada estudiante explique y revise una decisión. Una enseñanza humanista mantiene la precisión técnica, pero reconoce que aprender también implica construir autonomía y una relación respetuosa con los demás.

\needspace{5\baselineskip}\section*{Referencias}
{\setstretch{1.15}\setlength{\parskip}{6pt}
\par\hangindent=1.27cm\hangafter=1  Delors, J., et al. (1996). \textit{La educación encierra un tesoro}. UNESCO/Santillana.
\par\hangindent=1.27cm\hangafter=1  Scott, C. L. (2015b). \textit{El futuro del aprendizaje 2: ¿Qué tipo de aprendizaje se necesita en el siglo XXI?} (Investigación y prospectiva en educación, 14). UNESCO.
\par\hangindent=1.27cm\hangafter=1  Secretaría de Educación Pública. (2022). \textit{Plan de Estudio para la educación preescolar, primaria y secundaria}. SEP.
}

\end{document}
